\documentclass[a4paper]{article}
% Espanol (Mexico) con XeLaTeX: fontspec para fuentes Unicode, polyglossia
% para la division silabica y los nombres del documento en la variante mexicana.
\usepackage{fontspec}
\usepackage{polyglossia}
\setdefaultlanguage[variant=mexican]{spanish}
% polyglossia no traduce los nombres de listings.
\AtBeginDocument{\providecommand{\lstlistingname}{}\renewcommand{\lstlistingname}{Listado}%
  \providecommand{\lstlistlistingname}{}\renewcommand{\lstlistlistingname}{Índice de listados}}
\usepackage{amsmath,amssymb}
\usepackage{graphicx}
\usepackage[margin=2.35cm]{geometry}
\usepackage[most]{tcolorbox}
\usepackage{etoolbox}
\usepackage{listings}
\usepackage{xcolor}
\usepackage{hyperref}
\usepackage{booktabs,longtable,tabularx,array}
\usepackage{subcaption}
\usepackage{float}
\usepackage{enumitem}
\usepackage{microtype}
\usepackage{fancyhdr}
\usepackage{needspace}

\definecolor{kbBack}{HTML}{F2F6FA}
\definecolor{kbFrame}{HTML}{9DB4CC}
\definecolor{kbTitle}{HTML}{52708F}
\definecolor{ibBack}{HTML}{FBF5E7}
\definecolor{ibFrame}{HTML}{C9A961}
\definecolor{ibTitle}{HTML}{7D5F1E}
\definecolor{wbBack}{HTML}{FCF2F0}
\definecolor{wbFrame}{HTML}{D6A096}
\definecolor{wbTitle}{HTML}{9E4F43}
\definecolor{dbBack}{HTML}{F4F7F3}
\definecolor{dbFrame}{HTML}{AEC2AC}
\definecolor{dbTitle}{HTML}{5F7F64}

\tcbset{softbox/.style={
  enhanced,breakable,boxrule=0.8pt,arc=2pt,left=3mm,right=3mm,
  top=2mm,bottom=2mm,coltitle=white,fonttitle=\bfseries,
  toptitle=0.6mm,bottomtitle=0.6mm
}}
\newtcolorbox{knowledgebox}[1]{softbox,colback=kbBack,colframe=kbFrame,colbacktitle=kbTitle,title={#1}}
\newtcolorbox{importantbox}[1]{softbox,colback=ibBack,colframe=ibFrame,colbacktitle=ibTitle,title={#1}}
\newtcolorbox{warningbox}[1]{softbox,colback=wbBack,colframe=wbFrame,colbacktitle=wbTitle,title={#1}}
\newtcolorbox{dialoguebox}[1]{softbox,colback=dbBack,colframe=dbFrame,colbacktitle=dbTitle,title={#1}}

\lstset{
  language=Python,basicstyle=\ttfamily\small,
  keywordstyle=\color{kbTitle}\bfseries,stringstyle=\color{wbTitle},
  commentstyle=\color{dbTitle},breaklines=true,frame=single,numbers=left,
  numberstyle=\tiny\color{gray},captionpos=b,columns=fullflexible,
  keepspaces=true,showstringspaces=false,extendedchars=true
}
\BeforeBeginEnvironment{lstlisting}{\Needspace{12\baselineskip}}

\hypersetup{colorlinks=true,linkcolor=kbTitle,urlcolor=kbTitle,citecolor=wbTitle}
\setlist{itemsep=0.25em,topsep=0.4em}
\setlength{\parindent}{2em}
\setlength{\parskip}{0.25em}
\newcolumntype{L}[1]{>{\raggedright\arraybackslash}p{#1}}

\providecommand{\notetitle}{Notas del curso: [tema del curso]}
\providecommand{\noteauthors}{Elaboradas a partir de los materiales públicos de Stanford CS153}
\providecommand{\notedate}{Primavera de 2026}
\providecommand{\videochannel}{Stanford Online}
\providecommand{\videopublishdate}{2026}
\providecommand{\videoduration}{[duración del video]}
\providecommand{\videourl}{}
\providecommand{\videocoverpath}{cover.jpg}
\providecommand{\coursesite}{https://cs153.stanford.edu/}
\providecommand{\repourl}{https://github.com/hqhq1025/ai-course-notes}
\providecommand{\skillversion}{wdkns/wdkns-skills@39f1a04}

\newcommand{\figsource}[1]{\par\vspace{-0.3em}{\footnotesize\color{gray}\emph{Fuente: #1}}\par}
\newcommand{\teachervoice}[1]{\begin{knowledgebox}{Nota de clase}#1\end{knowledgebox}}
\newcommand{\term}[1]{\textbf{#1}}

\pagestyle{fancy}
\fancyhf{}
\fancyfoot[C]{\thepage}
\fancyfoot[R]{\footnotesize\color{gray}\href{\repourl}{AI Course Notes}}
\renewcommand{\headrulewidth}{0pt}
\renewcommand{\footrulewidth}{0pt}

\newcommand{\makecscover}{
\begin{titlepage}
\centering
\vspace*{0.7cm}
{\Large Stanford CS153 · Frontier Systems\par}
\vspace{0.8cm}
{\huge\bfseries \notetitle\par}
\vspace{0.6cm}
{\large \noteauthors\par}
\vspace{0.25cm}
{\large \notedate\par}
\vspace{0.9cm}
\ifdefempty{\videocoverpath}{}{
  \includegraphics[width=0.86\textwidth,height=0.43\textheight,keepaspectratio]{\videocoverpath}\par
}
\vfill
\begin{tcolorbox}[width=0.92\textwidth,colback=black!2!white,colframe=black!45,boxrule=0.7pt,arc=2pt]
\textbf{Autor o canal del video}: \videochannel\par
\textbf{Fecha de publicación}: \videopublishdate\par
\textbf{Duración del video}: \videoduration\par
\textbf{Sitio del curso}: \href{\coursesite}{\nolinkurl{\coursesite}}\par
\ifdefempty{\videourl}{}{\textbf{Enlace del video}: \href{\videourl}{\nolinkurl{\videourl}}\par}
\textbf{Normas de elaboración}: \skillversion\ + las normas source-first / teacher-voice / visual-QA de este repositorio
\end{tcolorbox}
\end{titlepage}
\tableofcontents
\newpage
}


\renewcommand{\notetitle}{Lecture 06: Infraestructura de IA a escala nacional: de la eficiencia energética y la difusión del cómputo a los Agent de gobierno}
\renewcommand{\noteauthors}{Elaborado a partir de la entrevista con Abdullah Alswaha, los subtítulos con marcas de tiempo y materiales oficiales/artículos académicos}
\renewcommand{\notedate}{Curso de invierno de 2025 · Versión reescrita en 2026}
\renewcommand{\videochannel}{Carga del curso histórico de Stanford CS153 (actualmente privado)}
\renewcommand{\videopublishdate}{2025-03-07}
\renewcommand{\videoduration}{52:17}
\renewcommand{\videourl}{https://www.youtube.com/watch?v=LriOr64E8D8}
\renewcommand{\videocoverpath}{cover.jpg}

\newcommand{\lecturefigure}[3]{%
\begin{figure}[H]
  \centering
  \includegraphics[width=0.96\textwidth]{images/#1}
  \caption{#2}
  \figsource{#3}
\end{figure}
}

\begin{document}
\makecscover

\section{Auditoría de fuentes: no es una lista de proyectos, sino una clase de diseño de sistemas a escala nacional}

Esta entrevista la impartió Abdullah Alswaha, ministro de Comunicaciones y Tecnologías de la Información de Arabia Saudita. Al verificarlo el 11 de agosto de 2026, el video público ya había pasado a ser private; el repositorio aún conserva 1,266 subtítulos con marca de tiempo y la portada oficial. Por ello, esta sesión reconstruye la argumentación de la clase a partir de los subtítulos y contrasta los mecanismos clave con Saudi Digital Government Authority, KFSHRC, Groq/Aramco Digital y artículos técnicos. En cuanto a las cifras de gigawatt, tamaño de mercado, costos, ahorros y productividad, si ningún documento de primera mano las da con el mismo criterio, este texto las presenta de manera uniforme como \textbf{estimación del ponente}, sin elevar lo dicho en el momento a hechos auditados.

\begin{importantbox}{La cadena causal de tres capas de toda la sesión}
\begin{enumerate}
  \item \textbf{Capa de hardware}: semiconductor, memory, interconnect y cooling determinan la useful compute per watt;
  \item \textbf{Capa de difusión}: la compute capacity solo se convierte en aplicaciones mediante la planificación de recursos, la inference de bajo costo y los workflow de cada sector;
  \item \textbf{Capa de gobernanza}: data, process, human accountability y public value determinan si un agent puede pasar de la demo a los servicios públicos.
\end{enumerate}
\end{importantbox}

\begin{warningbox}{Los tres tipos de evidencia deben separarse}
Las palabras textuales de la clase sirven para reconstruir motivaciones y juicios; los materiales oficiales, para confirmar organizaciones, proyectos o mecanismos vigentes; el análisis de ingeniería, para explicar costos, límites y conclusiones transferibles. Los tres pueden respaldarse entre sí, pero no sustituirse. En particular, la capacidad planeada, el impacto económico y los mercados futuros de los proyectos nacionales solo pueden enunciarse según su nivel de evidencia.
\end{warningbox}

\subsection{Technocrat y signal-to-noise ratio}

La clase describe a Alswaha como \term{technocrat} (tecnócrata): no porque «los ingenieros sean por naturaleza más aptos para gobernar», sino para subrayar que quienes dirigen infraestructura necesitan entender las restricciones de los sistemas, la medición, las fallas y la inversión a largo plazo. El ponente resume su método de juicio con la \term{signal-to-noise ratio} (SNR, relación señal-ruido). Si $P_s$ representa la fuerza de la evidencia relevante para el objetivo y $P_n$ la interferencia y la incertidumbre, puede escribirse
\[
\mathrm{SNR}=\frac{P_s}{P_n}.
\]
Aquí la fórmula es solo una analogía para la toma de decisiones: aumentar la SNR no significa ignorar las objeciones, sino definir primero la función objetivo, la calidad de los datos, la reversibilidad y el costo de fallar, y después extraer señales estables de la propaganda, las fluctuaciones de corto plazo y las cifras no verificadas.

\teachervoice{Alswaha describe su trayectoria profesional con learn--earn--return y considera «volver a Stanford» una oportunidad para transmitir problemas a la siguiente generación. Lo que de verdad vale la pena conservar del tono de la clase no es el relato de éxito personal, sino esto: el talento técnico debe dirigir su capacidad hacia problemas con valor público a largo plazo que, a la vez, puedan verificarse con ingeniería.}

\subsection{Resumen de la sección}

El núcleo de esta sesión no es evaluar la estrategia de un país, sino aprender cómo un sistema de IA a escala nacional se descompone en problemas verificables de hardware, aplicaciones y gobernanza. Todos los casos posteriores vuelven al mismo criterio: cuál es la entrada, dónde está el cuello de botella, cómo se mide el resultado y quién asume la falla.
\section{De la ola de internet a la ola de la IA: la infraestructura es solo una condición necesaria}

El ponente compara la IA actual con los inicios de internet: en la fase inicial, el capital y la atención se concentraron en routers, switches, accelerators y centros de datos; la etapa que realmente transformó la productividad llegó cuando telemedicine, education, GovTech y los procesos empresariales empezaron a reestructurarse. Por eso esta sección no pregunta «¿importa la capacidad de cómputo?», sino «¿cómo atraviesa la capacidad de cómputo las fronteras de una organización hasta convertirse en servicios estables y resultados medibles?».

\subsection{Cadena de difusión en cuatro pasos: capacity, application, organization, outcome}

Esta sección convierte la analogía histórica anterior en una cadena de sistema que puede verificarse. La difusión de la infraestructura puede dividirse en cuatro pasos: el primero aporta capacidad tecnológica; en el segundo aparece un application trigger concreto en un sector; en el tercero la organización modifica sus flujos de datos, permisos y formas de trabajo; solo en el cuarto se producen resultados en calidad del servicio, productividad o bienestar. Al leer la figura, identifique primero el owner, las entradas y los indicadores de cada paso; si cualquiera de los pasos se rompe, la inversión en hardware puede quedarse en baja utilización, proyectos de demostración o subsidios insostenibles.

\lecturefigure{01-infrastructure-diffusion.png}{La infraestructura solo produce resultados medibles cuando pasa por un detonante de aplicación y una reestructuración organizacional.}{Subtítulos locales 03:40--06:20; concepto redibujado.}

\begin{knowledgebox}{Lectura de la figura: primero busque los puntos de ruptura, no conecte directamente con el PIB}
Al leer la figura, pregunte de derecha a izquierda: ¿con qué indicador se mide el outcome objetivo? ¿Qué workflow cambia de verdad? ¿De qué datos, talento y adquisiciones depende la aplicación? Solo al final pregunte cuántos accelerators se necesitan. Si una propuesta solo describe la capacity del lado izquierdo, sin la adoption, el feedback y la retirement rule del lado derecho, sigue siendo un plan de recursos, no un plan de transformación.
\end{knowledgebox}

\begin{importantbox}{Expresión mínima de verificación de la difusión de infraestructura}
El valor de un proyecto puede descomponerse de forma aproximada como
\[
V \approx C\times U\times A\times Q,
\]
donde $C$ es la capacity disponible, $U$ es la tasa de utilización efectiva, $A$ es la adoption y $Q$ es la mejora de calidad o productividad que aporta cada unidad de uso. Si cualquiera de los términos se acerca a cero, ampliar solo $C$ difícilmente creará resultados. Esta expresión no sirve para una valuación precisa, sino para obligar a que la propuesta complete los eslabones intermedios.
\end{importantbox}

\begin{warningbox}{La analogía histórica no sustituye las restricciones actuales}
Tanto internet como la IA pasan por una etapa infrastructure-first, pero la IA está sujeta además a restricciones de datos de entrenamiento, actualización de modelos, costo de inferencia, seguridad del contenido, permisos de automatización y human accountability. El valor de la analogía está en recordar que «la difusión de las aplicaciones llega después del auge del hardware»; no permite concluir que ambas curvas industriales se repetirán al mismo ritmo o con los mismos ganadores.
\end{warningbox}

\subsection{Resumen de la sección}

La IA a escala nacional no puede medirse solo por el volumen de compras, los MW o los parámetros del modelo. El método transferible consiste en construir la cadena completa de capacity a outcome y establecer indicadores separados para utilization, adoption, quality y el mecanismo de salida.
\section{La pila de problemas de eficiencia energética: más allá de los FLOPS hay tres muros}

Al bajar de la difusión de las aplicaciones a la capa inferior, lo primero que aparece es la eficiencia energética. El ponente enumera cuatro frentes: semiconductor, memory, interconnect y cooling. Lo que en realidad quiere decir es que el throughput de un AI system está limitado por el recurso más lento y que los FLOPS pico son solo una cota superior. Esta sección unifica estas restricciones con un razonamiento de tipo roofline, para no llamar «faltan GPU» a todos los cuellos de botella.

\subsection{Von Neumann bottleneck y data movement}

El \term{von Neumann bottleneck} (cuello de botella de von Neumann) se refiere a que las unidades de ejecución del procesador están separadas del almacenamiento, por lo que los datos y las instrucciones deben moverse entre la memory hierarchy y el compute. En muchos AI kernels, la energía y el tiempo no se consumen en las multiplicaciones y sumas en sí, sino en leer weights, KV cache y activations, y en el intercambio entre dispositivos. Para una workload puede escribirse una cota inferior simplificada
\[
T_{step}\geq \max\left(\frac{F}{C},\frac{B_m}{BW_m},\frac{B_n}{BW_n}\right),
\]
donde $F$ es la cantidad de operaciones requerida y $C$ es el throughput de cómputo efectivo; $B_m$ y $BW_m$ son, respectivamente, el memory traffic y el memory bandwidth; $B_n$ y $BW_n$ son el network traffic y el interconnect bandwidth. La fórmula indica que, si solo se aumenta $C$, los otros dos términos todavía pueden dominar el tiempo.

\lecturefigure{02-efficiency-stack.png}{La eficiencia energética de la IA exige revisar a la vez semiconductor, memory, interconnect y facility.}{Subtítulos locales 06:20--12:00; redibujo conceptual.}

\begin{knowledgebox}{Lectura de la figura: las cuatro capas no son cuatro proyectos independientes}
Un nuevo semiconductor puede aumentar la rack density y, con ello, modificar el cooling y el suministro eléctrico; una \term{High Bandwidth Memory} (HBM, DRAM apilada de alto ancho de banda) más grande puede reducir la comunicación entre tarjetas, pero eleva el costo de encapsulado; el optical interconnect puede disminuir las pérdidas en la transmisión a larga distancia, pero introduce nuevos dispositivos y una nueva pila de software. La pregunta correcta es cuántos SLO-compliant tokens per watt obtiene la workload de extremo a extremo, no si el indicador pico de algún dispositivo luce mejor.
\end{knowledgebox}

\subsection{Qué resuelve y qué no resuelve el compound semiconductor}

Un \term{compound semiconductor} (semiconductor compuesto) está formado por dos o más elementos, por ejemplo GaN (gallium nitride, nitruro de galio) o SiC (silicon carbide, carburo de silicio). En alto voltaje, alta frecuencia y conversión de potencia pueden ofrecer mejor campo de ruptura o mejor eficiencia, por lo que son adecuados para fuentes de alimentación, radiofrecuencia y dispositivos optoelectrónicos; sin embargo, «migrar por completo desde el silicon» no es un interruptor único. Los chips lógicos, la memory, la power electronics y la photonics tienen requisitos distintos de materiales, rendimiento de fabricación y costo.

\begin{warningbox}{La analogía de dispositivos de la clase no debe tomarse como hoja de ruta de producto}
Los subtítulos mezclan switching voltage, leakage y compound materials en una sola respuesta rápida. Estas notas conservan la afirmación de que «la capa de dispositivos todavía tiene margen de eficiencia energética», pero no presentan ningún material como sustituto general de la GPU. La viabilidad técnica depende además de process integration, yield, packaging, toolchain y economías de escala.
\end{warningbox}

\subsection{Resumen de la sección}

La eficiencia energética de la infraestructura de IA es una propiedad del sistema de extremo a extremo. Si cualquiera de las capas (aritmética, memory, network o facility) se queda corta, un accelerator costoso puede quedar en espera; tanto la investigación como las compras deberían usar como indicador común el useful output de workloads reales.
\section{Memory Wall: el modelo crece más rápido que la capacidad de suministrar datos}

La sección anterior estableció las cotas inferiores de recursos; esta sección se centra en la memory wall, que se enfatizó repetidamente en clase. Gholami et al., en \emph{AI and Memory Wall}, analizan el desajuste entre el crecimiento del tamaño de los modelos, la memory capacity, el bandwidth y la comunicación. En un LLM, el training requiere parámetros, optimizer state, gradient y activation; la inference, además, debe gestionar la KV cache. Que el modelo «quepa» y que el modelo «reciba datos suficientes» son dos condiciones distintas.

\subsection{Muro de capacidad, muro de ancho de banda y muro de comunicación}

Esta sección descompone la «memory wall» en tres problemas medibles: si el modelo y su estado de ejecución caben, si pueden entregarse a compute por unidad de tiempo y si, al escalar a varias tarjetas, pueden sincronizarse de forma eficiente. Al leer la figura siguiente, no hay que agrupar el lado derecho bajo «falta memory», sino registrar por separado la evidencia de capacity, bandwidth y communication, porque cada tipo de problema requiere soluciones de hardware y software distintas.

\lecturefigure{03-memory-wall.png}{La esencia de la memory wall es el desajuste de velocidades entre compute throughput y data supply.}{Subtítulos locales 08:20--10:30; Gholami et al., \emph{AI and Memory Wall}.}

\begin{knowledgebox}{Lectura de la figura: primero, distinguir los tres tipos de wall}
La \textbf{capacity wall} pregunta si los parámetros, la KV cache y las activations caben en memory; la \textbf{bandwidth wall} pregunta si los datos pueden entregarse a compute por unidad de tiempo; la \textbf{communication wall} pregunta si la sincronización y la redistribución entre varios dispositivos absorben las ganancias del escalamiento. Las tres pueden presentarse al mismo tiempo, y resolver una de ellas expone el cuello de botella de la siguiente capa.
\end{knowledgebox}

\begin{importantbox}{La arithmetic intensity determina si el límite es «cómputo o ancho de banda»}
Se define la arithmetic intensity
\[
I=\frac{F}{B_m},
\]
es decir, cuántas operaciones se realizan por cada byte transferido. Si $I$ es muy baja, aunque el throughput pico del accelerator sea muy alto, es más probable que el kernel esté limitado por el memory bandwidth; el valor de fusion, tiling, quantization y cache reuse puede entenderse como aumentar la $I$ efectiva o reducir $B_m$.
\end{importantbox}

\subsection{Términos clave: SRAM, DRAM, HBM y storage}

Los subtítulos y los materiales comerciales suelen agrupar estos términos bajo «memory». El lector debe comprender primero el papel de cada uno en la hierarchy y, después, las diferencias entre Groq, NPU, GPU o las arquitecturas nuevas.

\begin{knowledgebox}{Concepto previo: cuatro niveles de almacenamiento}
\begin{tabularx}{\textwidth}{L{0.17\textwidth}L{0.24\textwidth}X}
\toprule
Término & Características principales & Papel y limitaciones en sistemas de IA \\
\midrule
SRAM & Static Random-Access Memory, caché de alta velocidad en el chip; baja latencia, pero costosa en área & Adecuada para datos pequeños y de uso frecuente, y para accesos deterministas; su capacidad es reducida y no equivale a un almacenamiento general para modelos. \\
DRAM & Dynamic Random-Access Memory; requiere refresco y tiene mayor capacidad & Memoria principal o memory externa del accelerator; su latencia y su ancho de banda dependen de la interfaz, el encapsulado y los accesos concurrentes. \\
HBM & High Bandwidth Memory, DRAM apilada de alto ancho de banda & Ofrece alto throughput mediante una interfaz ancha; la capacidad, el suministro de encapsulado, el costo y el consumo de energía siguen siendo restricciones. \\
SSD & Almacenamiento persistente, de gran capacidad pero con latencia mucho mayor que la memory & Guarda checkpoints, artifacts de modelos y datos fríos; normalmente requiere staging y no puede sustituir directamente a la HBM. \\
\bottomrule
\end{tabularx}
\end{knowledgebox}

\lecturefigure{04-memory-hierarchy.png}{Los distintos memory tiers intercambian latency, bandwidth, capacity y energy.}{Subtítulos locales 09:00--11:10; redibujo conceptual.}

\begin{knowledgebox}{Lectura de la figura: en el serving de modelos, los datos se ubican según su «temperatura»}
El kernel state más usado y los buffers pequeños permanecen en la SRAM on-chip; los weights activos y la KV cache compiten por la HBM; la host memory, más grande, se encarga del staging; el SSD guarda los modelos fríos y los checkpoints. Si el routing cambia de modelo con frecuencia, el arranque en frío y la transferencia de datos se convierten en el costo principal; por ello, el diseño del model portfolio debe coordinarse con la memory placement.
\end{knowledgebox}

\begin{warningbox}{«Tal chip solo sirve para texto» no es una conclusión estable}
En clase se usaron empresas concretas para ilustrar rápidamente las compensaciones entre SRAM, DRAM y HBM, pero las capacidades de los productos evolucionan con el compiler, el model support, la memory topology y el serving software. Estas notas conservan los principios de la hierarchy y no toman los juicios expresados en clase en 2025 como una clasificación permanente de productos.
\end{warningbox}

\subsection{Resumen de la sección}

La memory wall no es la escasez de un solo componente, sino un problema combinado de capacity, bandwidth, communication y software locality. Antes de elegir un chip, conviene cuantificar el working set del modelo, los bytes per token, el crecimiento de la KV cache y el traffic entre dispositivos.
\section{Economía de los centros de datos: de los MW de la grid a los useful tokens}

Tener chips no equivale a contar con capacidad de cómputo utilizable. La energía que entra de la red eléctrica pasa por el suministro y la distribución eléctrica, el UPS, el cooling, la network, el storage, la CPU y el accelerator, y al final solo una parte se convierte en useful work que cumple los objetivos de calidad y latencia. Esta sección convierte la discusión del ponente sobre gigawatts y facility cost en un marco de contabilidad de recursos, sin repetir cifras absolutas no auditadas.

\subsection{PUE, utilización y workload efficiency}

El \term{PUE} (Power Usage Effectiveness) se define como
\[
\mathrm{PUE}=\frac{P_{facility}}{P_{IT}},
\]
donde $P_{facility}$ es la potencia total del centro de datos y $P_{IT}$ es la IT load de los servidores, la network, el storage, etc. Cuanto más se acerca el PUE a 1, menor es el facility overhead; pero no indica si el accelerator se aprovecha plenamente ni si el modelo cumple el SLO al menor costo.

\lecturefigure{05-power-budget.png}{La facility power pasa por el overhead y la IT load antes de producir useful accelerator work.}{Subtítulos locales 10:30--14:50; esquema conceptual redibujado.}

\begin{knowledgebox}{Lectura de la figura: los tres indicadores de eficiencia son indispensables}
La \textbf{facility efficiency} se mide con el PUE; la \textbf{hardware utilization}, con el busy time del accelerator y los stalls de memoria y de network; la \textbf{workload efficiency}, con los tokens, samples o training progress obtenidos en relación con los recursos. Un clúster ocioso con un PUE bajo sigue siendo un desperdicio, y un sistema con alta utilización que genera resultados fallidos tampoco tiene valor.
\end{knowledgebox}

\begin{importantbox}{Useful compute per grid watt}
Puede usarse una expresión diagnóstica
\[
\eta_{grid}=\frac{U\cdot \eta_{workload}}{\mathrm{PUE}},
\]
donde $U$ es la tasa de utilización efectiva de los recursos de IT y $\eta_{workload}$ es la cantidad de trabajo válido producido por unidad de IT power. Aumentar $U$ puede depender del procesamiento por lotes, el scheduling y la demand aggregation; aumentar $\eta_{workload}$ puede depender de mejores kernels, model routing, quantization y cache.
\end{importantbox}

\subsection{Relojes de planificación: la electricidad, los edificios, los chips y los modelos no van sincronizados}

Lo más peligroso en los proyectos de escala nacional no es «pronosticar mal», sino amarrar entre sí decisiones de vida útil distinta. La grid y el site pueden planificarse a diez años; el facility shell requiere varios años de construcción; los accelerators y la rack density cambian con rapidez, y la model demand depende a su vez de los algoritmos y de los productos. El lector debe separar en capas la obra civil irreversible y los compute modules reemplazables.

\lecturefigure{06-planning-clocks.png}{La infraestructura nacional de IA opera al mismo tiempo en cuatro escalas temporales que entran en conflicto.}{Subtítulos locales 14:45--16:10; esquema conceptual redibujado.}

\begin{knowledgebox}{Lectura de la figura: gestionar el regret con modularity}
La capa de electricidad y terreno conserva headroom; la facility admite cooling de mayor densidad y retrofits de network; los racks y los accelerators se mantienen lo más modulares posible; el software conserva la portabilidad entre modelos y hardware. La modularity eleva el costo inicial, pero reduce la dependencia de trayectoria del tipo «cuando el centro de datos entra en operación ya no es adecuado para la nueva generación de racks».
\end{knowledgebox}

\begin{warningbox}{La capacidad planificada no es capacidad en línea}
Los 1.5--2 GW, los plazos de construcción y los costos unitarios mencionados en la sesión son cifras dadas en clase. Una evaluación formal debe distinguir entre announced, contracted, under construction, energized, installed, available y utilized capacity; de lo contrario, una misma cifra mezcla etapas distintas: el terreno, la electricidad, el building shell y los tokens que efectivamente pueden servirse.
\end{warningbox}

\teachervoice{La intuición de ingeniería del ponente es «llevar la electricidad hasta el accelerator en la mayor medida posible». Una versión más completa de los apuntes debería añadir: el accelerator solo es un activo productivo cuando lo usa una workload confiable; la capacity planning debe construir al mismo tiempo el scheduler, el software, el talento, el onboarding de modelos y el demand pipeline.}

\subsection{Resumen de la sección}

La economía de los centros de datos debe contabilizarse desde el grid watt hasta el useful outcome. El PUE, la utilization, la workload efficiency y los planning clocks determinan en conjunto el costo total; presentar cualquiera de estos indicadores por separado puede inducir a error.
\section{Tres AI frontier: generación, acción y control físico}

Los recursos de base se incorporan a distintas formas de aplicación. La clase divide el enfoque en generative, agentic y physical AI; el valor didáctico de esta clasificación es que el failure mode y la evidence de cada una son completamente distintos. El contenido generado puede evaluarse fuera de línea; un agent modifica el workflow state; un physical system afecta directamente a personas y equipos, por lo que los requisitos de validación, permisos y recuperación aumentan de forma escalonada.

\subsection{De la calidad de la salida a la seguridad de la acción}

La sección anterior explicó que los tres tipos de sistemas de IA tienen distintos radios de acción; esta sección convierte esa diferencia en requisitos de validación. Al leer la figura, observe el orden generative, agentic, physical: el sistema se amplía de «generar candidatos» a «modificar el estado digital», y luego a «controlar el mundo físico»; cada paso aumenta los permisos, el impacto, la demora en la detección y el costo de recuperación, por lo que no puede reutilizarse el mismo benchmark.

\lecturefigure{07-three-ai-frontiers.png}{Generative, agentic y physical AI requieren evidencia y rendición de cuentas progresivamente más estrictas.}{Subtítulos locales 16:10--20:20; redibujo conceptual.}

\begin{knowledgebox}{Lectura de la figura: un mismo benchmark no puede validar los tres niveles de sistema}
En generative AI pueden revisarse factuality, utility y diversity; en agentic AI también hay que revisar tool correctness, state transition, retry y exception handling; en physical AI es indispensable agregar sensor failure, control stability, safe stop, human override y la normativa aplicable en sitio. La puntuación del modelo es solo una parte de la evidencia del sistema.
\end{knowledgebox}

\begin{importantbox}{A mayor radio de acción, mayor costo de validación}
El riesgo puede entenderse de forma aproximada como
\[
R \propto P(failure)\times I(impact)\times D_{detect},
\]
donde $P(failure)$ es la probabilidad de falla, $I(impact)$ es el impacto y $D_{detect}$ es la demora en detectar y corregir. Cuando un Agent obtiene más permisos, aunque la tasa de error del modelo no cambie, el impacto y la demora de recuperación pueden amplificarse; por ello hay que agregar minimización de permisos, simulation, approval gate y rollback.
\end{importantbox}

\subsection{Salud: acelerar la búsqueda no significa que desaparezca la gobernanza}

La clase usa las formulaciones de fármacos y el trasplante de corazón robótico para ilustrar la entrada de la IA en la medicina. Los materiales oficiales del KFSHRC confirman que en 2024 se realizó un fully robotic heart transplant; aquí, «fully robotic» describe la ruta quirúrgica y el uso del robotic system, y no debe reescribirse como una autonomous surgery sin un médico responsable. De manera similar, la IA puede acortar la candidate generation y la formulation search, pero el clinical trial, la regulatory review y la evidencia de seguridad a largo plazo siguen existiendo.

\lecturefigure{08-healthcare-boundaries.png}{La IA médica puede comprimir la búsqueda en investigación, pero no puede omitir la evidencia clínica ni los límites de responsabilidad.}{Subtítulos locales 16:40--18:20; materiales oficiales del KFSHRC sobre el trasplante de corazón totalmente robótico.}

\begin{knowledgebox}{Lectura de la figura: no sumar las dos líneas de tiempo}
A la izquierda está el research loop: generar candidatos, predecir propiedades y programar experimentos; a la derecha, la care delivery: selección de pacientes, planificación quirúrgica, control médico y monitoreo posoperatorio. La IA puede mejorar algunos pasos de ambas cadenas, pero cada cadena tiene datos, ética, regulación y responsables distintos.
\end{knowledgebox}

\begin{warningbox}{«De diez años a dos» debe entenderse como un caso de clase}
Las cifras de duración que aparecen en los subtítulos no cuentan con material suficiente para confirmar su punto de inicio, su punto final ni el proyecto de enfermedad. Las notas conservan solo la conclusión transferible: la IA puede acortar la búsqueda de formulaciones y experimentos, pero eso no permite afirmar que el ciclo completo de drug development o de approval se comprima en la misma proporción.
\end{warningbox}

\subsection{Energía: el Agent debe integrarse al ciclo cerrado, no solo generar recomendaciones}

La industria energética dispone de sensor, maintenance, geology y operational history, y también tiene altos costos de seguridad. El ponente mencionó la corrosion y el drilling; el verdadero problema de sistema es cómo entra la recomendación del modelo en una work order, quién la aprueba, cómo regresan los resultados y cómo se detiene la operación ante una anomalía. Un agent sin outcome capture solo puede generar análisis y no logra un operational learning sostenible.

\lecturefigure{09-energy-agent-loop.png}{Un Agent vertical necesita que datos, recomendaciones, decisiones humanas y retroalimentación de resultados formen un ciclo cerrado.}{Subtítulos locales 18:20--20:20; redibujo conceptual.}

\begin{knowledgebox}{Lectura de la figura: por qué la domain data sigue sin ser suficiente}
Los datos históricos pueden reflejar equipos antiguos, estrategias anteriores y selection bias; la falta de datos de sensores puede producir false confidence; las operation exitosas a menudo carecen de etiquetas detalladas. Por ello, el Agent necesita uncertainty, operator rationale, counterfactual e incident review, en lugar de entregar la proprietary data directamente a un modelo de propósito general.
\end{knowledgebox}

\teachervoice{La clase considera las capacidades industriales existentes de un país como punto de partida de los AI use case: un país energético no necesita entrenar solo modelos «grandes y generalistas», sino que puede construir primero ventajas en workflow de alto valor, alta densidad de datos y con retroalimentación. Esto se acerca más a la ingeniería de sistemas que «primero tener el modelo y luego buscar el problema».}

\subsection{Resumen de la sección}

Las tres AI frontier comparten una base común de modelos y capacidad de cómputo, pero la diferencia real está en el radio de acción, el ciclo cerrado de datos y la estructura de responsabilidades. Cuanto más cerca se está de los servicios públicos y del mundo físico, menos puede un único benchmark representar la madurez de un sistema.
\section{Model Agnostic: sin religión de modelos, pero con disciplina de evaluación}

A partir de los casos de la industria, el ponente plantea la postura model agnostic: no construir una identidad en torno a open/closed, small/large ni a un vendor en particular, sino elegir el modelo según los objetivos de salvar vidas, eficiencia y productividad. Esta postura no sostiene que «todos los modelos son iguales»; exige establecer task contract, evaluation, routing, fallback y version governance más rigurosos.

\subsection{Portfolio selection en lugar de un winner único}

Esta sección convierte «no model religion» en un proceso de selección continuo, y no en un lema de neutralidad ante los proveedores. Al leer la figura, se parte del task contract y después se comparan los candidatos, la evidencia y la routing policy; así, aunque las capacidades y los precios de los modelos cambien con rapidez, el sistema puede explicar por qué cierto tipo de solicitud usa cierto modelo, cuándo se recurre al fallback, cómo se revierte el cambio y qué restricciones de gobernanza no pueden quedar anuladas por una mejora de rendimiento.

\lecturefigure{10-model-agnostic.png}{La estrategia model-agnostic tiene como eje la tarea, la combinación de candidatos, la evidencia y la routing policy.}{Subtítulos locales 20:20--23:40; figura conceptual redibujada.}

\begin{knowledgebox}{Lectura de la figura: cuatro pasos para evitar la model religion}
Primero se definen la calidad de la tarea, la latency, la privacy, el language y el cost; después se forman candidatos open/closed, small/large y general/specialist; luego se ejecutan offline/online eval y red-team; por último, la elección se escribe como una routing policy con capacidad de auditoría. Cuando llega un modelo nuevo, se vuelve a comparar, en lugar de sustituir automáticamente todas las rutas.
\end{knowledgebox}

\begin{importantbox}{La selección de modelos es una optimización con restricciones}
Puede escribirse como
\[
m^*=\arg\max_m Q(m)\quad
\text{s.t.}\quad L(m)\leq L_{max},\ C(m)\leq C_{max},\ G(m)=1,
\]
donde $Q$ es la calidad de la tarea, $L$ es la latency, $C$ es el costo y $G$ indica si se cumplen las restricciones de gobernanza, de datos y de despliegue. Si la distribución de tareas varía, la solución óptima puede ser routing o ensemble, y no un único modelo.
\end{importantbox}

\begin{warningbox}{Model agnostic no equivale a independencia de la cadena de suministro}
El modelo puede reemplazarse, pero el tokenizer, el tool schema, la safety policy, la observabilidad, los datos de fine-tuning y el evaluation history no necesariamente son portables. La verdadera portability requiere versioned interface, prompt contract, data export, fallback y rollback, no agregar un menú desplegable a la UI.
\end{warningbox}

\teachervoice{El ponente extiende «no technology religion» a «no model religion». La versión de ingeniería de esta frase es: mantener la arquitectura falsable, mantener al proveedor reemplazable y conservar evidencia comparable de cada migración.}

\subsection{Resumen de la sección}

Model agnostic es un portfolio guiado por evaluation, no la idea de que «cualquier modelo sirve». El contrato de la tarea, las restricciones, las versiones y el routing determinan juntos la elección; cuanto más a menudo se cambia de modelo, más se necesitan interfaces de sistema estables y pruebas de regresión.
\section{Difusión de la capacidad de cómputo: Capacity, Inference y brecha digital}

La combinación de modelos termina por apoyarse en la infraestructura. La clase calificó la inversión a gran escala como no-regret, pero desde el punto de vista sistémico ese juicio solo se sostiene si la capacity se aprovecha, la inference es asequible, las aplicaciones pueden entrar y los beneficios pueden retroalimentarse. Esta sección usa el flywheel para analizar cómo tiene éxito «comprar capacidad de cómputo» y explica también por qué puede ampliar, en lugar de reducir, la brecha digital.

\subsection{Utilization flywheel y token economics}

La sección anterior resolvió «qué modelo elegir»; esta responde «cómo lograr que la capacity a escala nacional se use de verdad». Al leer la figura hay que distinguir entre installed capacity, available service, utilization real y application demand: solo cuando mejoran a la vez la planificación, el onboarding de modelos, la confiabilidad y la adopción por parte de los usuarios puede bajar el unit cost y, a su vez, atraer más workload efectivo; añadir racks por sí solo no pone en marcha el volante.

\lecturefigure{11-infrastructure-flywheel.png}{La capacity solo forma una retroalimentación positiva a través de la utilization, el unit cost y la application demand.}{Subtítulos locales 27:10--32:20, 42:30--45:10; contexto oficial de Groq/Aramco Digital.}

\begin{knowledgebox}{Lectura de la figura: el token price más bajo no es el único objetivo}
Un precio bajo puede ampliar la experimentation, pero los usuarios también necesitan model availability, calidad en árabe y en dominios verticales, latency, data protection, reliability y support. Si el workload no es estable o el stack de software es difícil de usar, la capacity se convierte en un subsidio; si la demanda de aplicaciones crece y la power y la network no tienen margen, la promesa de precios bajos puede a su vez perjudicar el SLO.
\end{knowledgebox}

\begin{warningbox}{Hay que separar el momento de la clase de los anuncios posteriores}
La entrevista trata la colaboración de inferencia entre Groq y Aramco Digital tal como era visible en marzo de 2025; el anuncio posterior de ampliación en LEAP 2025 puede servir para ilustrar la evolución, pero no debe presentarse retroactivamente como capacidad que ya se había entregado en el momento de la clase. Por eso estas notas solo toman como conclusión del mecanismo «incorporar rápidamente múltiples modelos mediante onboarding y reducir la barrera de entrada a la inference».
\end{warningbox}

\subsection{Training e inference son dos formas de infraestructura}

El \term{training} (entrenamiento) suele actualizar parámetros mediante trabajos síncronos a gran escala y busca rendimiento y time-to-train; la \term{inference} (inferencia) usa un modelo ya entrenado para responder solicitudes y está restringida por latency, availability, procesamiento por lotes, KV cache, ubicación geográfica y costos continuos. Ambos comparten accelerators, pero no comparten exactamente la misma topology, el mismo scheduler ni la misma business demand.

\lecturefigure{12-training-inference.png}{El training es más concentrado y síncrono; la inference es más distribuida y está restringida por el SLO.}{Subtítulos locales 32:20--36:50; redibujo conceptual.}

\begin{knowledgebox}{Lectura de la figura: no fijar una proporción permanente}
El juicio de la clase sobre la proporción training/inference es una observación de un momento dado. El desarrollo de modelos, el fine-tuning, los synthetic data, el test-time compute, los agent loops y el crecimiento de usuarios cambian el workload mix. Un capacity plan debe usar un scenario range y conjuntos modulares, en lugar de extrapolar la proporción actual hasta el final de la vida útil de las instalaciones.
\end{knowledgebox}

\begin{importantbox}{La versión de recursos y la versión de capacidades de la brecha digital}
La brecha en su versión de recursos consiste en el access price, la network y la accelerator availability; la brecha en su versión de capacidades consiste en datos, talento, workflow redesign, evaluation y procurement. Reducir solo el token price puede mejorar la primera, pero no mejora automáticamente la segunda. Los programas públicos deben incluir el adoption support y el outcome measurement en el presupuesto de infraestructura.
\end{importantbox}

\subsection{Resumen de la sección}

La inversión en capacidad de cómputo solo se convierte en un productive asset cuando entra en el utilization flywheel. El training y la inference tienen formas de recursos distintas, y la inclusión digital no se reduce a tokens baratos; el software, las skills, los datos y la organizational adoption también forman parte de la infraestructura.
\section{Soberanía de datos y Data Embassy: ubicación, control y recuperación son problemas distintos}

Los sistemas de alcance nacional enfrentan la data sovereignty. En la clase se mencionaron la data embassy y la virtualización, pero este tipo de conceptos se simplifica con facilidad en «los datos están dentro o fuera del país». Esta sección trata solo la dimensión arquitectónica y no deriva de la entrevista conclusiones jurídicas concretas: location, control, jurisdiction, portability y continuity deben diseñarse por separado y verificarse en conjunto mediante contratos, controles técnicos y simulacros.

\subsection{Cinco preguntas en lugar de una location checkbox}

\term{data sovereignty} (soberanía de datos) describe a qué leyes nacionales, derechos de control y políticas están sujetos los datos; \term{data residency} (residencia de datos) es más acotada y se ocupa de la ubicación física o lógica donde se almacenan los datos; \term{data embassy} (embajada de datos) suele referirse a brindar protección soberana y continuidad a activos digitales críticos mediante acuerdos jurídicos especiales e infraestructura transfronteriza. No es un producto estándar de uso general; es necesario leer los tratados y textos legales específicos.

\lecturefigure{13-data-sovereignty.png}{La soberanía de datos exige diseñar a la vez location, control, jurisdiction y continuity.}{Subtítulos locales 32:20--36:50; antecedentes del material de Saudi DGA cloud-governance.}

\begin{knowledgebox}{Lectura de la figura: las cinco preguntas de la revisión de arquitectura}
¿Dónde se almacenan en reposo y dónde se procesan los datos? ¿Quién posee las keys y la privileged identity? ¿Qué jurisdicción puede exigir una disclosure? ¿Pueden los datos y los modelos someterse a un portable export? ¿Cómo se recupera el sistema ante una falla del provider, de la region o de la red nacional? Solo cuando las cinco preguntas tienen un owner y un test claramente definidos, la sovereignty deja de ser un eslogan y se convierte en un sistema de control.
\end{knowledgebox}

\begin{warningbox}{La ubicación no aporta por sí sola seguridad ni disponibilidad}
La localización en una sola región puede aumentar las correlated failures; la replicación en varias regiones, a su vez, puede introducir problemas transfronterizos y de consistencia. Encryption, key custody, identity, logging, supply-chain assurance, backup integrity y restore drill deben diseñarse junto con la location policy.
\end{warningbox}

\teachervoice{El valor de la clase está en tratar la resiliencia nacional como un infrastructure requirement y no como un cumplimiento normativo a posteriori. Estas notas agregan lo siguiente: cualquier propuesta de «embajada digital» debe traducir los compromisos jurídicos en key ownership, replication topology, failover authority y simulacros periódicos.}

\subsection{Resumen de la sección}

La data sovereignty es un problema de arquitectura multidimensional. La residency es solo uno de sus insumos; el control, la jurisdicción, la portability, la recuperación y la auditoría determinan si el sistema está realmente disponible en una crisis.
\section{AI × Space y nuevos mercados: de los sensores a las decisiones sobre recursos}

La discusión sobre space en la entrevista abarca communication, Earth observation, navigation, agricultura, agua y lunar robotics. Si solo se enumeran estos proyectos uno junto a otro, el resultado es un catálogo promocional; una estructura con mayor valor didáctico es la sensing-to-decision chain: la space infrastructure genera datos, la AI los interpreta, el gobierno o las empresas modifican la physical action y, después, los resultados se usan para calibrar el modelo.

\subsection{Communication, observation y harsh environment}

Esta sección reorganiza los proyectos de space vistos en clase como una sensing-to-decision pipeline. Al leer la figura, primero distingue las capacidades básicas que aportan la orbit y el sensor; después sigue cómo la imagery, el positioning y la connectivity llegan a la AI interpretation; por último, verifica si farming, water allocation o rover action reciben retroalimentación desde tierra. Solo cuando el último paso cierra el ciclo, los space data dejan de ser un producto informativo y se convierten en capacidad de gestión de recursos.

\lecturefigure{14-ai-space-stack.png}{La cadena de sistemas entre AI y space se extiende desde la infraestructura de percepción hasta las decisiones sobre recursos y las decisiones físicas.}{Subtítulos locales 24:00--27:10, 48:50--49:20; figura redibujada a partir del concepto.}

\begin{knowledgebox}{Lectura de la figura: la dificultad de la Earth observation está en el label y la intervention}
Las imágenes satelitales permiten identificar land use, vegetación o stress hídrico, pero la salida del modelo solo cambia los resultados si entra en un workflow de inspection, permit, irrigation o enforcement. Si la verdad de terreno es escasa, si la nubosidad y las variaciones estacionales son grandes y si las intervenciones de política pública no quedan registradas, se rompe la relación entre la precisión del modelo y el ahorro real de agua.
\end{knowledgebox}

\begin{warningbox}{Las proporciones de ahorro de agua y económicas mencionadas en clase son cifras del ponente}
Los subtítulos ofrecen varias proporciones sobre desertification, water saving, navigation y space-economy, pero carecen de un método y una línea base comunes. Este texto no toma esas cifras como conclusiones cuantitativas; solo conserva el patrón de sistema según el cual «geospatial sensing + AI + policy action» requiere una validación de ciclo cerrado.
\end{warningbox}

\subsection{NEOM y testbed: los nuevos mercados no eximen de la exigencia de evidencia}

Una ciudad nueva o un proyecto de gran escala puede reducir la legacy integration y ofrecer estándares unificados de identity, sensor e infrastructure, por lo que resulta adecuado como testbed; sin embargo, el carácter «greenfield» también amplifica los riesgos de scope, capital, governance y adoption. La aceptación del sistema debe partir igualmente de SLO de servicio en un ámbito reducido, unit economics, safety incidents y resident outcomes, en lugar de sustituir los datos de operación con la visión del proyecto.

\begin{importantbox}{Los cuatro niveles de evidencia de un testbed}
\begin{enumerate}
  \item el prototype funciona;
  \item el pilot es estable con una población y un entorno controlados;
  \item la production se sostiene con carga real, fallas y restricciones organizacionales;
  \item la transfer puede reproducirse en distintas ciudades, marcos regulatorios y condiciones de proveedores.
\end{enumerate}
Solo al alcanzar los dos últimos niveles queda demostrado que las capacidades de la plataforma son transferibles.
\end{importantbox}

\subsection{Resumen de la sección}

Lo esencial de AI × space no es la lista de satélites o rovers, sino si sensing, interpretation, decision y outcome forman un ciclo cerrado. Un greenfield testbed puede acelerar la experimentación, pero no reduce la exigencia de evidencia en producción ni de capacidad de transferencia.
\section{Nuevos paradigmas de cómputo: llevar el compute hacia los data}

En la segunda mitad de la clase se vuelve a la hardware research y se discuten in-memory, near-memory, photonics y compound semiconductor. Su objetivo común no es simplemente aumentar los peak FLOPS, sino reducir el data movement, la electrical/optical conversion y la synchronization. El lector debe ver estas direcciones como un workload-specific design space, y no como la lista definitiva de la próxima arquitectura de propósito general.

\subsection{In-memory y near-memory computing}

\term{in-memory computing} (cómputo en memoria) hace que el memory device ejecute parte de las operaciones en el mismo lugar donde almacena los datos; \term{near-memory computing} (cómputo cercano a la memoria) coloca el compute en el memory stack o junto a él, para acortar la ruta de transferencia. Ambos buscan mitigar el von Neumann data movement, pero enfrentan limitaciones de precision, device variation, endurance, programming model y workload mapping.

\lecturefigure{15-compute-near-memory.png}{El cómputo cercano a la memoria o en memoria reduce el movimiento de datos y, con ello, intercambia eficiencia energética por generalidad.}{Subtítulos locales 45:10--47:20; Sebastian et al. 2020.}

\begin{knowledgebox}{Lectura de la figura: por qué «mover menos datos» no equivale a ahorrar automáticamente la mitad de la energía}
La ganancia depende de si la operación puede mapearse a una memory primitive, de si la precisión de la conversión es suficiente, de si los ADC/DAC periféricos y el control son costosos, de si los datos deben moverse entre capas y de si el compiler puede aprovechar el hardware. Los porcentajes mencionados en clase deben entenderse como motivación de investigación, no como una garantía para cualquier workload.
\end{knowledgebox}

\subsection{Términos clave: Photonic, CXL y collectives}

\begin{knowledgebox}{Tres conceptos de interconexión que se confunden con facilidad}
\begin{tabularx}{\textwidth}{L{0.19\textwidth}L{0.25\textwidth}X}
\toprule
Término & Problema que resuelve & Mecanismo y límites \\
\midrule
Photonic interconnect & Reducir el consumo energético y la presión de densidad en transferencias de gran ancho de banda a larga distancia & Transmite datos mediante luz, pero la electro-optical conversion, el laser, el packaging y el control siguen teniendo un costo. \\
CXL & Coherent attachment / pooling de CPU, accelerator y memory & Mejora la combinación y el uso compartido de recursos; no equivale a eliminar la latency ni la contention de la memory remota. \\
Collectives & Comunicación colectiva entre varias GPU: all-reduce, all-gather, reduce-scatter, broadcast, entre otras & Es la semántica de software del distributed training/serving; el rendimiento depende de la topology, el message size, el overlap y la library. \\
\bottomrule
\end{tabularx}
\end{knowledgebox}

\begin{warningbox}{La architecture religion es tan peligrosa como la model religion}
Una arquitectura nueva puede ser muy potente en un kernel específico y, aun así, carecer de compiler, debugger, library, availability o multi-tenant isolation. Tanto en la adquisición como en la investigación conviene comparar el end-to-end workload, el migration cost y la ecosystem maturity, en lugar de basarse solo en demostraciones de dispositivos o en la eficiencia energética pico.
\end{warningbox}

\teachervoice{La invitación del ponente a los estudiantes fue directa: prestar atención al punto donde convergen las señales de publication, prototype y startup. Estas notas la reformulan como método de investigación: primero identificar el cuello de botella del sistema, después elegir un workload medible y, por último, confirmar que la ganancia del dispositivo se sostiene a través del compiler y del application stack.}

\subsection{Resumen de la sección}

Los nuevos paradigmas de cómputo giran en torno al data movement. In/near-memory, photonics, CXL y los nuevos materiales resuelven problemas de distintos niveles; la validación final sigue siendo la eficiencia energética, la corrección, la programabilidad y la confiabilidad de un workload de extremo a extremo.
\section{Robotics: se automatiza la tarea, no la responsabilidad}

Cuando el nuevo hardware y los nuevos modelos llegan a la physical AI, se enfrentan a la medicina, al robotaxi y a la harsh-environment robotics. La clase los plantea bajo una misma pregunta, pero su madurez de ingeniería es muy distinta. La pregunta que conecta esta sección es: qué task sustituye realmente el sistema, qué human role conserva y cómo sale de forma segura cuando ocurre una sensor failure o una control failure.

\subsection{Tres tipos de robotics con safety case distintos}

Un healthcare robot puede mejorar la precision, pero el surgeon y el hospital asumen la clinical responsibility; el robotaxi percibe y controla de manera continua en vías abiertas, por lo que necesita un operational design domain y soporte remoto; el lunar/water-scarcity rover opera con retrasos de comunicación y en entornos extremos, y prioriza la autonomy, la fault tolerance y la energy management. Los tres no pueden compartir un mismo «robotics benchmark» genérico.

\begin{importantbox}{El safety case mínimo de la physical AI}
\begin{enumerate}
  \item Definir el operational design domain y las condiciones de prohibición;
  \item Indicar la sensor coverage, la uncertainty y el degradation mode;
  \item Explicar el human override, el safe stop y la recovery;
  \item Registrar la action, el state, la versión de software/model y los incident;
  \item Validar con escenarios reales y con rare event, no solo con la tasa de éxito promedio.
\end{enumerate}
\end{importantbox}

\begin{warningbox}{La transferencia de trabajo no es un beneficio social automático}
Delegar las repetitive task a un robot puede liberar mano de obra, pero también puede provocar desajustes de habilidades, desigualdad entre regiones y más trabajo de supervisión. La propuesta del sistema debe diseñar al mismo tiempo el retraining, el transition support, la liability y el service access; no puede tratar el «trabajo de mayor valor» como un resultado que ocurre por sí solo.
\end{warningbox}

\subsection{Resumen de la sección}

La madurez de la robotics depende del task boundary, el safety case, el human role y la incident recovery. Automatizar una acción no equivale a transferir la responsabilidad, y mucho menos a que toda una profesión o industria ya esté automatizada.
\section{Government Agent: de la Task Demo al Workflow System}

La última pregunta de toda la clase es la más cercana a la agent engineering actual. Alswaha distingue entre la AI enhancement del legacy B2B SaaS y el AI-native player, y sostiene que la dificultad no está solo en el modelo, sino en los data, el business model y el human-in-the-loop. Dicho con más precisión, el task-level success solo se convierte en un sistema de servicio público si cruza tres puertas: workflow integration, multi-agent coordination y public accountability.

\subsection{Task, workflow y multi-agent: tres expansiones del estado}

Un \term{task-level agent} completa drafting, classification o lookup sobre entradas locales; un \term{workflow-level agent} hace avanzar un case a través de varios sistemas, permisos y estados de larga duración; un \term{multi-agent system} hace que varios agents se repartan el trabajo, compartan memory, coordinen herramientas y resuelvan conflictos. Cada expansión aumenta el state space, la failure interaction y la audit complexity.

\lecturefigure{16-government-agent-maturity.png}{Un Agent gubernamental debe cruzar, en orden, task, workflow, multi-agent y public accountability.}{Subtítulos locales 49:50--52:17; redibujo conceptual.}

\begin{knowledgebox}{Lectura de la figura: por qué la task demo es la que con más facilidad tiene éxito}
En una task demo, las entradas pueden filtrarse manualmente, el estado es corto, los permisos son pocos y las fallas pueden ignorarse; un workflow real se topa con campos faltantes, cases duplicados, ownership entre departamentos, deadlines, appeals y policy exceptions. Multi-agent introduce además shared memory, conflicting plans, deadlock y cascading tool error; por ello, la tasa de éxito en producción no puede extrapolarse a partir de un benchmark de una sola ronda.
\end{knowledgebox}

\subsection{Data layer, human-in-the-loop y business model}

\term{human-in-the-loop} (humano en el ciclo) no significa «que una persona confirme con un clic cualquier resultado», sino turnar los casos de alto riesgo, los de baja confianza, las excepciones de política y las vías de apelación a una persona con la autoridad correspondiente, y recopilar sus razones como evidencia para mejorar. La data layer requiere canonical entity, lineage, quality rule, permission y retention; el business model, por su parte, responde quién paga el despliegue, los errores, la revisión humana y la actualización continua.

\begin{importantbox}{Matriz de aceptación de un Agent de servicio público}
\begin{tabularx}{\textwidth}{L{0.18\textwidth}L{0.29\textwidth}X}
\toprule
Dimensión & Pregunta clave & Evidencia medible \\
\midrule
Calidad de la tarea & ¿La salida es correcta, completa y explicable? & expert review, error taxonomy, calibration \\
Integridad del proceso & ¿Hace avanzar correctamente el state y las exception? & end-to-end case success, rework, deadline miss \\
Equidad y derechos & ¿Se puede apelar? ¿Produce un impacto diferenciado entre grupos? & subgroup audit, appeal outcome, override reason \\
Seguridad y privacidad & ¿Los permisos, los datos y las tool están minimizados? & access log, red-team, leakage/abuse test \\
Viabilidad económica & ¿La automatización reduce de verdad el total cost / cycle time? & compute + integration + human review + incident cost \\
\bottomrule
\end{tabularx}
\end{importantbox}

\begin{warningbox}{«El modelo no es la dificultad» tampoco debe tomarse como absoluto}
Los data y el workflow suelen ser el principal cuello de botella, pero la language coverage, el reasoning, la tool reliability, la calibration y la robustness del base model siguen limitando el sistema. La conclusión más precisa es esta: la capacidad del modelo es una condición necesaria, y el éxito en producción lo determina el producto de model, data, process, people, policy y economics.
\end{warningbox}

\teachervoice{La teacher voice más transferible de toda la clase aparece al final: algunos equipos conectan los vertical data directamente al modelo y solo obtienen task augmentation; un sistema verdaderamente AI-native rediseña la data creation, el human feedback y la workflow ownership.}

\subsection{Resumen de la sección}

El límite de un Government Agent no está en la interfaz de chat, sino en el case state, los permisos, los datos, las excepciones, las apelaciones y la responsabilidad. Una task demo puede estar impulsada por el modelo; un workflow system debe estar impulsado por una arquitectura sociotécnica completa.
\section{Síntesis y ampliación}

\subsection{Conclusiones centrales}

\begin{importantbox}{Doce conclusiones transferibles}
\begin{enumerate}
  \item La ola de la IA solo genera productividad si atraviesa la difusión de aplicaciones y la reestructuración organizacional.
  \item El useful compute está sujeto a las restricciones conjuntas de semiconductor, memory, interconnect y facility.
  \item El memory wall comprende tres desajustes: capacity, bandwidth y communication.
  \item SRAM, DRAM, HBM y SSD deben usarse por niveles según la temperatura del working-set.
  \item El PUE solo mide el facility overhead; la utilization y la workload efficiency son igual de importantes.
  \item Grid, facility, accelerator y model demand avanzan a ritmos distintos; la modularity sirve para gestionar el regret.
  \item La IA generative, la agentic y la physical requieren evidence y accountability distintas.
  \item Acelerar la búsqueda médica no permite omitir la clinical governance, y la cirugía robótica tampoco equivale a autonomous surgery.
  \item Model agnostic significa un portfolio guiado por evaluation (evaluation-driven), no la supresión de estándares.
  \item La capacity solo es un productive asset si entra en el flywheel de utilization--cost--adoption.
  \item La data sovereignty debe diseñar al mismo tiempo location, control, jurisdiction, portability y recovery.
  \item La dificultad decisiva de un Government Agent está en el workflow, los data, el human-in-the-loop y la responsabilidad económica.
\end{enumerate}
\end{importantbox}

\subsection{Tarea práctica: diseñar un sistema nacional de AI capacity-to-outcome}

Elige uno de estos ámbitos: healthcare, energy, education o public administration, y completa el siguiente diseño:

\begin{enumerate}
  \item Traza la cadena de contabilidad de recursos desde el grid power hasta el useful outcome, y define PUE, utilization y workload efficiency;
  \item Estima el working set, el memory traffic, el network traffic y el latency SLO, y determina si el cuello de botella está en compute, memory o network;
  \item Define tres demand scenarios para training e inference, así como la modularización de la instalación y la estrategia de adquisiciones;
  \item Establece una evaluation model-agnostic y una routing policy que incluyan fallback y version rollback;
  \item Diseña los controles de las cinco dimensiones de la data sovereignty y plantea simulacros de falla de region/provider;
  \item Extiende el agent de la task al workflow, y enumera state, permission, exception, appeal y human gate;
  \item Define indicadores de adoption, service quality, equity y total cost, y explica cuándo detener o reducir el proyecto.
\end{enumerate}

\begin{knowledgebox}{Criterios de aceptación}
Una propuesta de alta calidad no se limita a escribir «comprar GPU, desplegar un modelo grande, construir una plataforma». Debe separar las decisiones irreversibles de power/facility del hardware/model reemplazable, separar las estimaciones del ponente de los hechos verificados y proponer variables intermedias medibles que vayan de la capacity al public outcome.
\end{knowledgebox}

\subsection{Lecturas adicionales}

\begin{enumerate}
  \item Gholami et al., \href{https://arxiv.org/abs/2403.14123}{AI and Memory Wall}.
  \item Sebastian et al., \href{https://www.nature.com/articles/s41565-020-0655-z}{Memory devices and applications for in-memory computing}.
  \item Saudi Digital Government Authority, \href{https://dga.gov.sa/sites/default/files/2024-08/Cloud%20Computing%20and%20its%20Role%20in%20Accelerating%20Digital%20Transformation%20and%20its%20Sustainable%20Impact%20in%20the%20Government%20Sector%20-%20V1.0.pdf}{Cloud Computing and Digital Transformation}.
  \item KFSHRC, \href{https://www.kfshrc.edu.sa/en/news/2024/09/kfshrc-pioneers-worlds-first-fully-robotic-heart-transplant}{World's first fully robotic heart transplant}.
  \item Groq, \href{https://groq.com/news_press/groq-and-aramco-digital-announce-partnership-to-establish-worlds-largest-inferenc-data-center-in-saudi-arabia}{Aramco Digital inference partnership} y \href{https://groq.com/news_press/saudi-arabia-announces-1-5b-expansion-to-fuel-ai-powered-economy-with-groq}{later LEAP 2025 expansion}.
  \item Saudi Data and AI Authority, \href{https://sdaia.gov.sa/en/SDAIA/about/Pages/Strategy.aspx}{National Strategy for Data \& AI}.
\end{enumerate}

\end{document}
